% !TeX root = ../papers/nc-p3.tex
\section{曲線グラフにおけるただ一つの単純な純関係と直線への還元}
\label{sec:curve-graph-single-pure}
\begin{theorem}\label{thm:curve-graph-single-pure}
四辺の全射影ファイバーが空または一点で、点グラフが一次元の標準AS正則窓では、
各辺の関係空間とSegre多様体のスキーム論的交叉は空か被約な一点である。
もし $R_0$ に $u\otimes v$ があれば、共通点グラフの全ての一次元被約成分は全位置で直線となり、
位置0では全て $u^\perp$ の平面に、位置1では全て $v^\perp$ の平面に入る。
\end{theorem}
\begin{proof}[長さ三の無限小純関係]
$R_0$ の純関係の長さ3の部分スキームがあれば、その線型包 $P$ はベクトル次元3である。
二次元以下なら長さ3のスキームが乗る射影直線上で全二次小行列が恒等零となり、
純関係の直線が生じて前定理に反する。両文字の射影の包は各々高々3で、
前定理の部分支持上界から各々丁度3である。第一の包を消す文字 $p$ を取ると $R_0(p)$ は
残る三関係の像であり、点ファイバー条件により階数3である。従って $p$ は点グラフ上にある。
その前点を $y$ とする。$w(y)$ の末尾縮約が張る三関係の空間は $R_0$ の
$p$ で消える部分に入る。両者とも三次元で後者は $P$ だから等しい。
従って $w(y)$ は三つの支持がいずれも三次元で、非零の一文字縮約に純テンソルがある。
定理\ref{thm:ternary-point-subalgebras}により非正則となる。以上は任意の辺に使える。
\end{proof}
\begin{proof}[長さ二の純関係と四次評価鉛筆]
長さ2の部分スキームの包を $P$ とする。次元は2、両支持も2である。
$P$ の補空間を $E\subset R_0$ とし、完全対合でテンソルを分解したときのこの補空間の部分を
$B\in E\otimes V_2\otimes V_3$ と書く。$P$ の第一支持の消滅直線 $L$ 上に
四行四列の一次評価 $F(p):E\to V_1$ ができ、全点で階数3以上である。
$B$ は位置2,3とも全支持である。一文字で $B$ が消えれば、残る $P$ の二項だけからなる
三テンソルは少なくとも一方の端支持が2以下となり、一点ファイバー条件に反する。

まず $\det F$ が恒等零とする。各 $p\in L$ に前点 $y(p)$ と次点 $q(p)$ があり、
次数の一致によりどちらの射も次数1である。$\ker F(p)$ の射を $\alpha:L\to\mathbb P E$ とする。
$4\mathcal O(-1)\to4\mathcal O$ の核は $\mathcal O(-1-\deg\alpha)$、
余核は $\mathcal O(\deg q)$ であり、$\deg\alpha+\deg q=3$、従って $\deg\alpha=2$。
隣接点での零関係から $B(y(p))$ は $\alpha(p)\otimes d(p)$ であり、全支持のため非零である。
$y(p)$ からの基点なし線型写像としては次数1だが、このSegre射の次数は2以上であり矛盾する。

$\det F\ne0$ ならその零因子 $T$ は長さ4で、$L$ と点グラフのスキーム論的交叉に等しい。
実際単元な三次小行列の近くで核点を消去すれば残るのは行列式のみである。
核線 $\alpha:T\to\mathbb P E$ は完全な四次元切断空間を使う。
転置した鉛筆の列 $0\to4\mathcal O(-1)\to4\mathcal O\to\mathcal L_T\to0$ で
余核は $T$ 上可逆、$H^0(4\mathcal O)\to H^0(\mathcal L_T)$ が同型だからである。
点グラフから $y:T\to X_3$ を戻し、$D=H^0(\mathcal O_T)$ とおく。
$B(y)=\alpha\otimes d$ はここでも定階数1の因数分解である。
実際 $F(p)B(y)=0$ は点グラフ上の恒等式であり、$\ker F$ は $T$ 上の直線束である。
線束を自明化し、$d$ 上でいたるところ非零となる一文字を選ぶと、
$B(y)$ のその文字による縮約は $\alpha$ の単元倍である。
従って $y$ の線型座標も $D$ 全体を張る。$d$ の線型座標も同様に $D$ 全体を張る。
さもなくば四文字の一つで $B(y)$ が $T$ 上恒等零となり、$y$ の全線型包を用いると
$B$ の位置2の全支持に反するからである。

これら三つの同定の下で、$B$ により $R_2$ に入る四空間は $D$ の全ての乗法行列の空間である。
実際 $E^*$ と位置2の文字を縮約し、位置3の文字を $y$ 上で読むと対応する二関数の積が戻る。
$D$ は直線上の因子のFrobenius環であり、完全積対合を取ればこの四空間は乗法の双対像となる。
その純関係スキームは $T$ 自身の閉埋込みを含む。各試験環への点評価は積を保ち、
商の二因子の評価の単純積を与える。用いる線型系が $D$ 全体なのでこれは閉埋込みである。
この中の長さ3の部分因子に先の排除を適用すれば矛盾する。
これで長さ2もなく、前定理の有限性からスキーム論的な主張を得る。
重根での固有直線と乗法双対の明示検算は\path{computations/verify_curve_graph_pure_pair.py}にある。
\end{proof}
\begin{proof}[一純点が全曲線に課す条件]
$P=k(u\otimes v)$ とし、今度は $E$ を五次元補空間とする。
ある整曲線成分の位置0で $p(u)=0$ と仮定する。その正規化上の前点を $y$、次点を $q$ とし、
各文字束の共通次数を $d>0$ とする。$F(p):E\otimes\mathcal O\to V_1\otimes L_0$ は階数3一定で、
余核は $L_0L_1$、核 $K$ は階数2、次数 $-2d$ である。
$B(y):V_2^*\otimes\mathcal O\to K\otimes L_3$ は全射である。
実際 $w(y)$ の末尾縮約像は三次元で、その $E$ への射影で失う次元は高々1、像は $K$ に入る。
標的は大域生成で行列式の次数が零なので、この商を与えるGrassmann写像は定数となる。
従って $B(y)$ の $V_2$ 側の支持は一つの固定二平面 $D_2$ である。

$D_2$ を消す非零点 $x$ を取ると $B(x)$ の位置3の支持は $y$ の曲線の消滅空間に入る。
$y$ が三次元以上のベクトル空間を張れば、この支持は一次元以下となる。
$P$ の一項を戻しても $w(x)$ の位置3の端支持が2以下となり矛盾する。
従って $y$ の像は直線、$d=1$ で全位置とも直線である。
次に位置2の対応点は $B(y)$ の位置2の全支持を消すので、位置2の直線は
$\mathbb P(D_2^\perp)$ そのものである。
これで位置2,3の二直線の\emph{直積全体}で $B$ は消える。
その直積への $R_2$ の制限は $(1,1)$ のグラフの式一本で、階数2の双線型形式である。
実際グラフは射影同型のグラフであり、全式が零なら一点ファイバーでなくなる。
従って制限四テンソルはその形式と $u\otimes v$ の積である。
位置2または3の各点で相手側の直線に沿う縮約が非零となるので、
位置1の対応点は $v$ を消し、位置0の対応点は $u$ を消す。
最初に $q(v)=0$ から始めても語順の反転で同じ結論になる。
各整曲線上では $p(u)q(v)=0$ なので少なくとも一方から始められる。
\end{proof}
